\documentclass{article}
\usepackage{graphicx} % Required for inserting images
\usepackage{amssymb}
\usepackage{amsmath}
\usepackage[margin=1in]{geometry}
\usepackage{listings}
\usepackage{float}

\lstset{
    language=Matlab,
    basicstyle=\ttfamily\small,
    keywordstyle=\color{blue},
    commentstyle=\color{green!60!black},
    stringstyle=\color{red},
    numbers=left,
    numberstyle=\tiny,
    stepnumber=1,
    numbersep=5pt,
    frame=single,
    breaklines=true,
    showstringspaces=false
}
    

\title{Numerical Analysis of Variable Separability in a Boundary-Value Diffusion Equation in Polar Coordinates [Report]}
\author{Kevin Bedoya}

\begin{document}

\maketitle

\section{Methodology}

\subsection{Variable Separability Analysis}

To determine whether the computed diffusion density exhibits approximate variable separability, the numerical solution was analyzed using matrix-based techniques motivated by the structure of a separable function.

Let

\[
D \in \mathbb{R}^{(m-1)\times n}
\]

denote the truncated density matrix obtained from the numerical solution. The final radial ring was excluded from the analysis because it corresponds to the absorbing boundary condition

\[
\phi(r=1,\theta)=0,
\]

and therefore consists entirely of zeros. Each row of \(D\) corresponds to a discrete radial location, while each column corresponds to a discrete angular location.

If the diffusion density is separable,

\[
\phi(r,\theta)=R(r)\Theta(\theta),
\]

then the sampled density matrix admits the factorization

\[
D = \mathbf{r}\mathbf{\theta}^{T},
\]

where \(\mathbf r\) contains samples of the radial component and \(\mathbf \theta\) contains samples of the angular component. Consequently, a perfectly separable density matrix has rank one.

To investigate this property numerically, three complementary diagnostics were employed.

\subsubsection{Angular Trajectory Alignment}

The angular trajectory matrix was defined as

\[
A = D^{T}.
\]

Each column of \(A\) contains the angular density profile

\[
\phi(\theta \mid r=r_i)
\]

at a fixed discrete radius \(r_i\).

For a separable solution,

\[
\phi(\theta \mid r=r_i)
=
R(r_i)\Theta(\theta),
\]

so every angular trajectory differs only by a multiplicative constant. Therefore all columns of \(A\) should be mutually aligned.

Each column of \(A\) was normalized to unit Euclidean norm,

\[
\widehat{a}_i
=
\frac{a_i}{\|a_i\|_2},
\]

and the normalized matrix

\[
\widehat{A}
=
\begin{bmatrix}
\widehat{a}_1 &
\widehat{a}_2 &
\cdots &
\widehat{a}_{m-1}
\end{bmatrix}
\]

was formed. The cosine similarity matrix

\[
C_A
=
\widehat{A}^{T}\widehat{A}
\]

was then computed. The entries of \(C_A\) represent pairwise cosine similarities between normalized angular trajectories.

Two metrics were used to quantify alignment:

\[
E_{\mathrm{mean}}
=
\frac{
\left\|
\,|C_A|-1
\right\|_F
}
{\sqrt{N}},
\]

where \(N\) denotes the total number of entries in \(C_A\), and

\[
E_{\mathrm{max}}
=
\max_{i,j}
\left|
\,|C_{A,ij}|-1
\right|.
\]

The first quantity measures the average deviation from perfect alignment, while the second measures the worst-case deviation. Values near zero indicate that the angular trajectories are nearly proportional across all radii.

\subsubsection{Radial Trajectory Alignment}

An analogous procedure was performed using the radial trajectory matrix

\[
B=D.
\]

Each column of \(B\) contains the radial density profile

\[
\phi(r \mid \theta=\theta_j)
\]

at a fixed angular position \(\theta_j\).

For a separable solution,

\[
\phi(r \mid \theta=\theta_j)
=
\Theta(\theta_j)R(r),
\]

so all radial trajectories should also be proportional. After column normalization, a cosine similarity matrix

\[
C_B
=
\widehat{B}^{T}\widehat{B}
\]

was formed and the same mean and maximum alignment errors were computed. Values close to zero indicate that the radial trajectories are nearly proportional across all angular positions.

\subsubsection{Rank-One Dominance}

In addition to trajectory alignment, the full density matrix was examined using the singular value decomposition

\[
D = U\Sigma V^{T}.
\]

Let

\[
\sigma_1 \ge \sigma_2 \ge \cdots
\]

denote the singular values of \(D\). The rank-one energy ratio was defined as

\[
\eta
=
\frac{\sigma_1^2}
{\sum_{k}\sigma_k^2}.
\]

This quantity measures the fraction of the total matrix energy captured by the dominant singular mode. A perfectly separable solution has rank one after discretization. 

\[
\eta = 1.
\]

\section{Results}

\begin{figure}[H]
    \centering
    \includegraphics[width=0.8\linewidth]{48x48_T=2_angular.png}
    \label{fig:placeholder}
\end{figure}


\begin{figure}[H]
    \centering
    \includegraphics[width=0.8\linewidth]{48x48_T=2_radial.png}
    \label{fig:placeholder}
\end{figure}



\begin{lstlisting}

Frobenius Cosine Similarity Metric (phi v. theta):  8.772448094240826e-05
Max Deviation (phi v. theta):  8.772448094240826e-05
Frobenius Cosine Similarity Metric (phi v. radius):  0.02092182614888484
Max Deviation (phi v. radius):  0.028431069706455703
Rank-1 dominance ratio (general):  0.9999839882679707
    
\end{lstlisting}



\end{document}

